% rhythm: article display math: a full last line into an unnumbered display and out of it, then the same around a numbered equation
% boundary: par-display Birch span
% boundary: display-par Cedar span
% boundary: par-display Dogwood span
% boundary: display-par Elm span
\documentclass{article}
\usepackage{amsmath}
\usepackage{fontspec}
\setmainfont{OpenSans-Regular.ttf}[Path=../corpus/fonts/, BoldFont=OpenSans-Bold.ttf, ItalicFont=OpenSans-Italic.ttf, BoldItalicFont=OpenSans-BoldItalic.ttf]
\usepackage{unicode-math}
\setmathfont{FiraMath-Regular.otf}[Path=../corpus/fonts/]
\begin{document}
Alder sets a paragraph whose last line runs the full measure, so that the
display below it opens with the long skip and never the short one that a
short last line would earn: this sentence is long enough to fill its lines.
\[ a + b = \text{Birch} \]
Cedar follows the unnumbered display, and it too runs on long enough that its
last line is full before the numbered equation below it takes its own skip,
which is the long one again because the line before it is not short at all.
\begin{equation} c + d = \text{Dogwood} \end{equation}
Elm closes the page after the numbered equation.
\end{document}
